% Propietario conservado: /Users/ruben/Documents/New project/output/EXTREMA_MEDIA_RAZON_RECIPROCIDAD_ES_EN_20260918/ARTICULO/ES/source/libro/reciprocidad_recuperacion.tex
\chapter{Coordonnée électronique, échelle opératoire et réversibilité des lectures}
\label{x_reciprocidad:rec:recuperacion}

\section{La structure algébrique dans l’expression électronique complète}

L’expression de la section~\ref{x_reciprocidad:sec:electron} se compose maintenant avec le récupérateur radial, en maintenant ses facteurs et sa réalisation dimensionnelle. Le préfacteur électronique provient des projecteurs conditionnels de somme et de produit sur \(\ell^2(\{1,\ldots,9\}^3)\). Dans le plan principal correspondant,
\[
 P=\begin{pmatrix}1&0\\0&0\end{pmatrix},\qquad
 Q=\begin{pmatrix}1/4&\sqrt3/4\\\sqrt3/4&3/4\end{pmatrix},
 \qquad\|[P,Q]\|=\frac{\sqrt3}{4}.
\]
L'égalité se vérifie car \([P,Q]=(\sqrt3/4)\left(\begin{smallmatrix}0&1\\-1&0\end{smallmatrix}\right)\).
De plus,
\[
 S=2P-I,\qquad J=\frac4{\sqrt3}[P,Q],\qquad
 S^2=I,\quad J^2=-I,\quad SJ=-JS.
\]
C’est pourquoi \(n_\pm=(S\pm J)/2\) vérifient \(n_\pm^2=0\). Le bloc conserve des générateurs hyperbolique, elliptique et nilpotents dans la réalisation algébrique déjà dérivée. Le coefficient \(\sqrt3/4\) et la normalisation \(4/\sqrt3\) remplissent des fonctions distinctes.

Nous définissons, avec les coordonnées et les entiers régionaux préalablement construits,
\begin{align}
 \Delta_4&=\frac{\pi-e\log\pi}{270}\left(1+\frac\pi{729}\right),\nonumber\\
 B_e&=\frac{\sqrt3}{4}
       \left[\exp\!\left(\frac\varphi{\pi^2}\right)-22\alpha^3\right]
       (1+15\Delta_4),\qquad
 \kappa_e=80-54\alpha+6\Delta_4.
 \label{x_reciprocidad:rec:eq:electron-basal}
\end{align}
Ainsi, \(\varepsilon_e^{\beta}=B_eR_{\rm act}^{\kappa_e}\). La section préalable d'action et son retour fixent
\[
 \delta_{\rm ret}=H_5-\eta_{\rm ret}=\frac{90}{\pi}\mathcal C_\pi,
 \qquad R_{\rm act}=1+\frac{\delta_{\rm ret}}{\eta_{\rm ret}}.
\]

\begin{corollary}[Réalisation radiale de la coordonnée électronique]
\label{x_reciprocidad:rec:cor:electron}
Dans la branche \(\alpha>0\), \(\eta_{\rm ret}>0\), \(|\xi|<1\), on a
\begin{align}
 \varepsilon_e^{\beta}={}&
 \frac{\sqrt3}{4}
 \left[\exp\!\left(\frac\varphi{\pi^2}\right)-22\alpha^3\right]
 (1+15\Delta_4)\nonumber\\
 &\times\left[1+\frac{90\mathcal C_\pi}{\pi\alpha}
 \frac{1+R_\star^4}{499-501R_\star^4}\right]^{80-54\alpha+6\Delta_4}.
 \label{x_reciprocidad:rec:eq:electron-radio}
\end{align}
Dans la coordonnée de Klein de la carte marquée,
\begin{equation}
 R_{\rm act}=1+
 \frac{\delta_{\rm ret}}{\alpha(-500x_K-1)}.
 \label{x_reciprocidad:rec:eq:accion-klein}
\end{equation}
Les domaines correspondants sont
\[
 \frac1{500}<\xi<1,\qquad
 0<R_\star^4<\frac{499}{501},\qquad
 -1<x_K<-\frac1{500}.
\]
\end{corollary}
\begin{proof}
On substitue \eqref{x_reciprocidad:rec:eq:eta-radio} dans \(R_{\rm act}=1+\delta_{\rm ret}/\eta_{\rm ret}\) puis dans l’expression électronique complète \eqref{x_reciprocidad:eq:el-formula-completa}. La carte de Klein donne \(\xi=-x_K\). Les intervalles s’obtiennent à partir de \(\eta_{\rm ret}=\alpha(500\xi-1)>0\) et des correspondances monotones déjà démontrées.

\end{proof}

La substitution conserve le déficit régional, la normalisation torsionnelle, le caractère de retour et la norme d'incompatibilité. Elle exprime la même coordonnée au moyen d'un lecteur géométrique récupérable. Le rayon et \(\alpha\) restent corrélés par leur génération ; ils ne constituent pas des paramètres indépendants ajustés à une masse. La réalisation dimensionnelle reste
\[
 E_e^{\beta}=\varepsilon_e^{\beta}\mathcal U_E,\qquad
 m_e^{\beta}=E_e^{\beta}/c_{\rm int}^2.
\]
La décennie d’action déjà incorporée dans les sections de Planck n’est pas multipliée de nouveau sur une coordonnée exprimée dans la carte d’énergie.

Une seconde forme exacte, utile pour isoler la correction, est
\[
 \widehat\delta_{\rm ret}=\frac{90\mathcal C_\pi}{\pi H_5},\qquad
 R_{\rm act}=(1-\widehat\delta_{\rm ret})^{-1},\qquad
 \varepsilon_e^{\beta}=B_e(1-\widehat\delta_{\rm ret})^{-\kappa_e}.
\]
La positivité des deux sections garantit \(1-\widehat\delta_{\rm ret}>0\). Les trois expressions sont des identités du même retour, non des prédictions indépendantes.

\section{Séparation de l’échelle et de la forme dans une fibre positive de masses}

La loi opératoire de retour reçoit une fibre positive finie de routes avec ses lecteurs. Pour formuler sa conséquence matricielle sans dépendre de noms d’espèces, soit \(\mathcal H\) de dimension \(d\ge1\), soit \(B=B^*\) l’opérateur des exposants et soit \(M^0\) défini positif. Le retour est la congruence
\begin{equation}
 M^\beta=R_{\rm act}^{B/2}M^0R_{\rm act}^{B/2}.
 \label{x_reciprocidad:rec:eq:masa-operador}
\end{equation}
Dans la base de routes du lecteur diagonal, \(Be_\gamma=\kappa_\gamma e_\gamma\) et
\(M^0e_\gamma=m_e^{(0)}\chi_{\rm ref}(\sigma_\gamma-\sigma_e,\eta_\gamma-\eta_e)e_\gamma\). Les signatures et le caractère sont ceux de la construction des routes ; le théorème suivant s’applique à chaque fibre positive sans substituer son sélecteur.

\begin{theorem}[Décomposition uniforme et unimodulaire du retour]
\label{x_reciprocidad:rec:thm:masa}
Si \(R_{\rm act}>0\), définissons
\[
 \bar\kappa=\frac{\tr B}{d},\qquad B_0=B-\bar\kappa I,
 \qquad S_0=R_{\rm act}^{B_0/2}.
\]
Alors
\begin{align}
 M^\beta&=R_{\rm act}^{\bar\kappa}S_0M^0S_0,
 \qquad\det S_0=1,\label{x_reciprocidad:rec:eq:masa-forma}\\
 \det M^\beta&=R_{\rm act}^{\tr B}\det M^0.\label{x_reciprocidad:rec:eq:masa-det}
\end{align}
Avec \(\widehat M=M/(\det M)^{1/d}\),
\begin{equation}
 \widehat M^\beta=S_0\widehat M^0S_0.
 \label{x_reciprocidad:rec:eq:masa-normalizada}
\end{equation}
\end{theorem}
\begin{proof}
La partie scalaire de \(B\) commute avec \(B_0\). Par calcul fonctionnel, \(R_{\rm act}^{B/2}=R_{\rm act}^{\bar\kappa/2}S_0\). Dans une base orthonormale de vecteurs propres de \(B\),
\[
 \det S_0=R_{\rm act}^{\frac12\sum_\gamma(\kappa_\gamma-\bar\kappa)}=1.
\]
La multiplication des déterminants prouve \eqref{x_reciprocidad:rec:eq:masa-det}. En divisant \eqref{x_reciprocidad:rec:eq:masa-forma} par la racine \(d\)-ième de ce déterminant, le facteur scalaire se simplifie. La preuve ne requiert pas la commutation de \(B\) avec \(M^0\) ; elle requiert en revanche le caractère défini positif pour normaliser par le déterminant.
\end{proof}

Dans la base diagonale conjointe du propriétaire, les quotients individuels sont
\[
 q_\gamma=\frac{m_\gamma^\beta}{m_\gamma^{(0)}}R_{\rm act}^{-\bar\kappa}
 =R_{\rm act}^{\kappa_\gamma-\bar\kappa},\qquad
 \prod_\gamma q_\gamma=1.
\]
L’inversion formelle \(R_{\rm act}\mapsto R_{\rm act}^{-1}\) inverse chaque \(q_\gamma\). Elle est distincte de la réciprocité radiale orientée, qui conserve \(R_{\rm act}\) par \eqref{x_reciprocidad:rec:eq:eta-orientada}. La congruence conserve le déterminant normalisé, et non les valeurs propres individuelles ou la somme des masses. Le secteur nul est exclu de cette normalisation positive. Pour \(d=1\), \(B_0=0\) : le retour de l’électron réside intégralement dans le facteur uniforme et reste dans \eqref{x_reciprocidad:rec:eq:electron-radio}.

La substitution de \eqref{x_reciprocidad:rec:eq:accion-klein} agit également sur \eqref{x_reciprocidad:rec:eq:masa-operador} par calcul fonctionnel. Si
\[
 q(\xi)=1+\frac{\delta_{\rm ret}}{\alpha(500\xi-1)}>0,
\]
alors
\[
 M^\beta=e^{\frac12(\log q(\xi))B}\,M^0\,
          e^{\frac12(\log q(\xi))B}.
\]
La généalogie de chaque exposant et de chaque multisection est conservée. La structure commune avec le flot de \(K\) est l’élimination d’une composante uniforme suivie d’une déformation de déterminant un ; on n’identifie ni \(B_0\) à \(\operatorname{diag}(u)\), ni \(R_{\rm act}\) à \(\rho\).

\section{Compression circulaire et identification du transport}

\begin{theorem}[Récupération bilatérale d'un opérateur unitaire]
\label{x_reciprocidad:rec:thm:operador}
Soit \(C\) unitaire dans un espace de Hilbert et soient
\[
 T=\frac{8I+C}{9},\qquad D=\frac{\sqrt8(C-I)}9.
\]
Alors
\begin{equation}
 T^*T+D^*D=I,\qquad
 v=Tv-\frac{Dv}{\sqrt8},\qquad Cv=Tv+\sqrt8Dv.
 \label{x_reciprocidad:rec:eq:unitario}
\end{equation}
\end{theorem}
\begin{proof}
Le développement donne
\[
 81T^*T=65I+8(C+C^*),\qquad
 81D^*D=16I-8(C+C^*).
\]
La somme est \(81I\). Les autres identités résultent de combinaisons linéaires des deux définitions. Pour une valeur propre circulaire \(|z|=1\), le bilan est \(1-|(8+z)/9|^2=(8/81)|1-z|^2\).
\end{proof}

La paire archivée récupère aussi bien l’état que son image. Pour le repère inversible \(O_K=(K,SK,\ldots,S^{11}K)\) déjà construit, les douze réponses forment
\[
 Y=(CK,CSK,\ldots,CS^{11}K),\qquad C=YO_K^{-1}.
\]
Si \(CS=SC\), une seule réponse vectorielle \(y=CK\) détermine toutes les colonnes : \(Y=(y,Sy,\ldots,S^{11}y)\). Dans le cas général, douze essais vectoriels sont nécessaires. La récupération ne réduit pas cette information à un unique scalaire.

Le générateur réel diagonal \(H_K=\operatorname{diag}(h_i)\) admet la carte unitaire
\[
 C_K=(H_K+iI/2)(H_K-iI/2)^{-1}.
\]
Chaque valeur propre réelle se transforme en un nombre de module un. De plus,
\[
 C_K-I=i(H_K-iI/2)^{-1},\qquad
 H_K=\frac i2(C_K+I)(C_K-I)^{-1}.
\]
Les deux inverses existent, car \(H_K\) est autoadjoint et son spectre est réel. L'archive de \eqref{x_reciprocidad:rec:eq:unitario} récupère donc cette réalisation finie de la direction de \(K\). Sa dimension douze et son domaine demeurent distincts de ceux d'un opérateur spectral de zêta.

\section{Hiérarchie des invariants et portée démonstrative}

La complétion \(1000=729+271\), de lecture normalisée \(1=729/1000+271/1000\), maintient sa fonction générative. Les conservations ultérieures s’organisent selon l’objet transporté :
\begin{center}
\small
\begin{tabularx}{\linewidth}{@{}p{.27\linewidth}Xp{.25\linewidth}@{}}
\toprule
Réalisation & Invariant & Récupération ou réciprocité\\
\midrule
Archivo bilateral & Norme conjointe de lecture et de mémoire & État et image par l’opérateur\\
Ellipse d’action & Produit \(a_Sb_S=2\hbar\) & Rayon et orientation angulaire\\
Flux de \(K\) & Produit de douze rayons & Inversion \(t\mapsto-t\)\\
Réseau marqué déformé & Covolume et isométrie ternaire & \(\Lambda_t^*=\Lambda_{-t}\)\\
Lecture thêta--Mellin & Pôles et résidus & Différence entière et équation fonctionnelle\\
Fibre positive des masses & Déterminant après retrait de l’échelle commune & Congruence par \(S_0\)\\
\bottomrule
\end{tabularx}
\end{center}

La conservation du volume n'implique pas la conservation de la norme de chaque vecteur. La conservation polaire permet de faire varier la partie entière. La récupération de l'action exige de conserver l'orientation. Ces distinctions expriment le contenu opératoire de la conservation évolutive : chaque transformation maintient certaines données, en modifie d'autres et conserve le registre qui permet de reconstruire la correspondance.

La composition étend la fonction de \(K\) déjà démontrée dans l’article. Le registre participe à la lecture de \(\alpha\) ; sa direction se récupère à partir de la mémoire ; cette direction détermine une géodésique radiale et une déformation réticulaire réciproque ; la famille réticulaire possède une lecture spectrale stable. La section d’action fournit le lien avec la coordonnée électronique et avec la congruence positive de l’opérateur de masse. On obtient ainsi une chaîne d’applications aux domaines, inverses et invariants explicites.


Le supplément distingue les preuves écrites des preuves formalisées. \texttt{FlujoReciproco.lean} contient dix théorèmes sur le flux diagonal, le centrage, les déterminants, l’inversion et le relèvement à vingt-quatre dimensions. \texttt{RadioKlein.lean} contient onze théorèmes rationnels sur les cartes, les bornes, les réciprocités et la récupération orientée. Les preuves hyperboliques, les substitutions complètes, la séparation de l’opérateur des masses et l’argument thêta--Mellin figurent dans le texte ; leur formalisation n’est pas attribuée à la compilation de ces deux fichiers. Les reçus documentaires enregistrent la provenance et la reproduction, avec une portée distincte de celle d’une démonstration mathématique.
